\documentclass[12pt]{article}
\usepackage[utf8]{inputenc}
\usepackage[T1]{fontenc}
\usepackage{hyperref}
\title{Hattic Language Resources for K-12 Curriculum Development}
\author{Compiled by Claude Sonnet 5.5 (claude-sonnet-5-5) for Andr\'as Kornai}
\date{}
\begin{document}
\maketitle

\section{Introduction}
Hattic (ISO 639-3 xht; Glottocode hatt1246) was a non-Indo-European language of ancient central and northern Anatolia. It is generally treated as a language isolate. Hattic speakers left no written records of their own, and the language was dead by the Hittite New Empire (c.\ 1400--1190 BCE). It is known only from about 150 short passages in Hittite cuneiform tablets. These are religious texts: rituals, incantations, litanies and myths. There are no native speakers and no known living community or school programme. Any curriculum would be a heritage and cultural-awareness programme built on a small, specialist-only evidence base.

\section{Primary Educational Resources}
I found no primers, textbooks or school materials for Hattic in my searches. Nothing suitable for K-12 pupils was found. Teachers would have to adapt scholarly material.

\section{Linguistic and Reference Materials}
\begin{itemize}
\item Wikipedia, ``Hattians'' (history and culture of the people): \url{https://en.wikipedia.org/wiki/Hattians}
\item Britannica, ``Hattian language'' (encyclopedic overview): \url{https://britannica.com/topic/Hattian-language}
\item Palaeolexicon, ``The Hattians and the Hattic language'' (word-list and overview page): \url{https://palaeolexicon.com/Hattic}
\item Goedegebuure, ``Hattian Alignment'' (scholarly paper on Hattic grammar): \url{https://voices.uchicago.edu/goedegebuure/files/2020/05/Goedegebuure-Hattian-Alignment.pdf}
\item Academia.edu listing of papers on Hattic: \url{https://www.academia.edu/Documents/in/Hattic_Language}
\end{itemize}
Search summaries also mention the Catalogue des textes hittites (CTH), which catalogues the published Hattic texts. They also mention earlier studies by Laroche, Kammenhuber, Girbal, Klinger, Schuster, Soysal and Taracha. I saw no URLs for these, so they are not in the JSON list.

\section{Religious and Community Resources}
The surviving Hattic material is itself religious, for example the myth of the Moon God who fell from the Sky, which exists in Hattic and Hittite versions. I found no living religious or community organisation for Hattic.

\section{Digital and Audio Resources}
I found no audio, video or software resources for Hattic. The language has no native speakers, and any pronunciation would be reconstruction.

\section{Curriculum Development Context}
Very little exists. A realistic K-12 unit would use the Hattian people, their history and the surviving myths as cultural content in English. A few reconstructed words or phrases could be used, labelled clearly as uncertain. Ancient-history, linguistics and world-languages classes are the most plausible places for it.

\end{document}